\documentclass[../main.tex]{subfiles}
\ifSubfilesClassLoaded{\setcounter{chapter}{4}}{}
\begin{document}

\chapter{Struktur Kontrol Perulangan}

\begin{subcpmk}
  \item Sub-CPMK 2.2: Menyusun program dengan perulangan \texttt{for} dan \texttt{while} termasuk pola bercabang
\end{subcpmk}

\SesiRPS{5}{2.2}{$\sim$170 menit}{CPMK-2}

\section{Sarana dan prasarana}

Python 3.10+, editor, terminal.

\section{Prosedur kerja}

\begin{enumerate}[leftmargin=*]
  \item Praktik \texttt{for}, \texttt{while}, \texttt{break}/\texttt{continue} dari contoh.
  \item Kerjakan latihan loop bersarang sesuai aktivitas.
  \item Dokumentasikan pola iterasi pada laporan singkat.
\end{enumerate}

\section{Perulangan \texttt{for}}

\texttt{for} digunakan untuk iterasi pada rentang atau koleksi. Tidak seperti beberapa bahasa lain, \texttt{for} di Python mengiterasi \textit{item} dari sequence secara langsung.

\begin{lstlisting}[caption=for dengan range dan iterasi langsung (Tutorial Python ss4.2)]
# Iterasi pada range
for i in range(5):
    print(i, end=" ")
# 0 1 2 3 4

# Iterasi langsung pada string
for huruf in "Python":
    print(huruf, end="-")
# P-y-t-h-o-n-

# Iterasi pada list
buah = ["apel", "jeruk", "mangga"]
for b in buah:
    print(b)
\end{lstlisting}

\begin{konsep}
Saat melakukan iterasi dan perlu memodifikasi koleksi, lebih aman untuk bekerja pada \textit{salinan}: \texttt{for b in buah[:]:}. Tutorial Python resmi (§4.2) juga menyarankan menggunakan fungsi seperti \texttt{enumerate()} untuk mendapatkan indeks sekaligus nilai (dibahas di modul 6).
\end{konsep}

\section{Fungsi \texttt{range()}}

\texttt{range(a, b)} menghasilkan $a, a+1, \ldots, b-1$. \texttt{range(a, b, step)} menyertakan langkah (Tutorial Python §4.3):

\begin{lstlisting}[caption=Variasi range]
list(range(5))          # [0, 1, 2, 3, 4]
list(range(2, 8))       # [2, 3, 4, 5, 6, 7]
list(range(0, 20, 5))   # [0, 5, 10, 15]
list(range(10, 0, -2))  # [10, 8, 6, 4, 2]  -- menurun
\end{lstlisting}

\section{Perulangan \texttt{while}}

Cocok bila jumlah iterasi tidak pasti:

\begin{lstlisting}[caption=while]
n = int(input("n: "))
i = 1
while i <= n:
    print(i * i)
    i += 1
\end{lstlisting}

\section{\texttt{break} dan \texttt{continue}}

\texttt{break} menghentikan seluruh loop; \texttt{continue} melompat ke iterasi berikutnya:

\begin{lstlisting}[caption=break dan continue]
# continue: cetak bilangan ganjil saja
for x in range(10):
    if x % 2 == 0:
        continue
    print(x, end=" ")
# 1 3 5 7 9

# break: berhenti saat menemukan bilangan negatif
angka = [3, 7, -1, 5, 2]
for a in angka:
    if a < 0:
        print(f"\nDitemukan angka negatif: {a}, berhenti.")
        break
    print(a, end=" ")
\end{lstlisting}

\section{Klausa \texttt{else} pada Loop}

Ini adalah fitur \textbf{unik} Python yang sering mengejutkan programmer dari bahasa lain. Blok \keyword{else} pada loop dieksekusi \textbf{setelah loop selesai secara normal} (tidak dihentikan oleh \texttt{break}). Sangat berguna untuk algoritme pencarian.

\begin{lstlisting}[caption=else pada loop for (Tutorial Python ss4.5)]
# Mencari bilangan prima (contoh dari tutorial resmi)
for n in range(2, 10):
    for x in range(2, n):
        if n % x == 0:
            print(n, "=", x, "*", n//x)
            break
    else:
        # Loop tidak menemukan faktor -> n adalah prima
        print(n, "adalah bilangan prima")
\end{lstlisting}

\begin{lstlisting}[caption=else pada loop while: pencarian nilai]
daftar = [5, 3, 8, 1, 9, 2]
target = 7
i = 0
while i < len(daftar):
    if daftar[i] == target:
        print(f"Ditemukan {target} di indeks {i}")
        break
    i += 1
else:
    # Keluar dari while tanpa break -> target tidak ditemukan
    print(f"{target} tidak ada dalam daftar")
\end{lstlisting}

\begin{konsep}
Pola \texttt{for/while ... else} adalah idiom Python yang elegan untuk pencarian: jika loop selesai tanpa \texttt{break}, berarti kondisi pencarian tidak terpenuhi dan blok \texttt{else} dieksekusi.
\end{konsep}

\section{\texttt{pass} dalam Loop}

Seperti pada percabangan, \texttt{pass} dapat digunakan dalam loop sebagai placeholder:

\begin{lstlisting}[caption=pass dalam loop]
# Loop yang belum diimplementasikan isinya
while True:
    pass  # Placeholder - akan diisi dengan logika nanti
\end{lstlisting}

\section{Loop Bersarang}

Digunakan untuk pola matriks sederhana, pola bintang, atau kombinatorika kecil:

\begin{lstlisting}[caption=Pola bintang dengan loop bersarang]
n = 5
for baris in range(1, n + 1):
    print("*" * baris)
# *
# **
# ***
# ****
# *****
\end{lstlisting}

\begin{aktivitas}
  \item Cetak bilangan 1--100 yang habis dibagi 7 menggunakan \texttt{for}.
  \item Baca bilangan $n$, cetak faktorial $n!$ dengan \texttt{while} (asumsikan $n$ kecil).
  \item Gambar pola segitiga bintang tinggi $n$ dari masukan.
  \item Gabungkan loop dan \texttt{if}: jumlahkan bilangan positif sampai pengguna memasukkan 0 sebagai penanda (gunakan \texttt{while}).
  \item Implementasikan pencarian nilai dalam list menggunakan \texttt{for...else}: cetak pesan khusus bila nilai tidak ditemukan.
  \item Buat program yang mencari semua bilangan prima antara 2 dan 50 menggunakan pola \texttt{for...for...else}.
\end{aktivitas}

\begin{latihan}
  \item Kapan memilih \texttt{for} dibanding \texttt{while}?
  \item Apa beda \texttt{break} dan \texttt{continue}? Beri contoh masing-masing satu baris dalam loop.
  \item Jelaskan dengan kata-kata Anda sendiri apa yang dilakukan klausa \texttt{else} pada \texttt{for}. Di situasi apa ini berguna?
  \item Jelaskan kompleksitas waktu perkiraan untuk loop bersarang dua tingkat dengan \texttt{range(n)} pada keduanya.
  \item Refleksi: kesalahan \textit{off-by-one} sering terjadi pada \texttt{range}; bagaimana Anda menghindarinya?
\end{latihan}

\begin{asesmen}
\textbf{Latihan terbatas waktu}

\begin{enumerate}
  \item Loop manakah yang mencetak tepat tiga angka?
  \begin{enumerate}
    \item \texttt{for i in range(3): print(i)}
    \item \texttt{for i in range(1, 3): print(i)}
    \item \texttt{for i in range(4): print(i)}
    \item \texttt{for i in range(0, 10, 3): print(i)}
  \end{enumerate}
  \item Blok \texttt{else} pada loop \texttt{for} dieksekusi ketika:
  \begin{enumerate}
    \item Loop dijalankan setidaknya sekali
    \item Loop berakhir karena \texttt{break}
    \item Loop selesai tanpa \texttt{break}
    \item Kondisi loop bernilai \texttt{False} sejak awal
  \end{enumerate}
\end{enumerate}

\textbf{Tugas}: program yang membaca $n$ dan menghitung jumlah kuadrat bilangan 1 sampai $n$ dengan \texttt{for}. Tambahkan program \texttt{cari\_prima.py} yang mencetak semua bilangan prima di bawah 100 menggunakan pola \texttt{for...else}.
\end{asesmen}

\begin{checklist}
  \item Saya dapat memakai \texttt{range} dengan satu, dua, atau tiga argumen
  \item Saya dapat menulis \texttt{while} yang berhenti dengan kondisi yang benar
  \item Saya memahami \texttt{break} dan \texttt{continue}
  \item Saya memahami dan dapat menggunakan klausa \texttt{else} pada \texttt{for} dan \texttt{while}
  \item Saya dapat menyelesaikan masalah sederhana dengan loop bersarang
\end{checklist}

\begin{rangkuman}
Perulangan mengotomatisasi tugas berulang. \texttt{for} cocok untuk iterasi terkendali; \texttt{while} untuk kondisi umum. Klausa \texttt{else} pada loop adalah fitur unik Python yang berguna untuk algoritme pencarian --- blok \texttt{else} dieksekusi hanya jika loop selesai tanpa \texttt{break}.

\textbf{Referensi:} {\small Python Tutorial §4.2--4.6 \url{https://docs.python.org/3/tutorial/controlflow.html}}

\textbf{Kata kunci:} \code{for}, \code{while}, \code{range}, \code{break}, \code{continue}, \code{else}, \code{pass}
\end{rangkuman}

\section{Referensi sesi}

\begin{itemize}[leftmargin=*]
  \item \textbf{[16]} W3Schools, \textit{Python Tutorial}.
  \item \textbf{[17]} Programiz, \textit{Python Programming Tutorial}.
  \item \textbf{[8, 9]} \textit{The Python Tutorial}, perulangan.
\end{itemize}

\end{document}
